\section{Conclusions}
\label{sec:conclusions}

This paper reexamines the prevalent approach to performing system health management and makes the case for why keeping system health management functions separate from decision making functions can be inefficient and/or ineffective. We also present the case for why prognostics is only meaningful in a limited set of circumstances and, even then, needs to be driven by decision making requirements.

We then explain the rationale for unifying (not just integrating) system health management with decision making and outline an approach for accomplishing that. We also propose keeping emergency response functionality separate to guarantee timely response, provide dissimilar redundancy, and allow for offline computation, validation, and verification of emergency response policies.

We believe that the proposed unification approach will improve the performance of both decision making and system health management, while potentially simplifying informational architectures, reducing sensor suite sizes, and combining modeling efforts. The approach is scalable with respect to system complexity and the types of uncertainties present. It also puts a variety of existing and emerging computational methods at the disposal of system designers.

